% !TeX root = ../main.tex

% Toggle: set to true for final bound version,
% false for submission version (reviewers on separate sheet)
\newif\ifreviewersinbinding
\reviewersinbindingtrue  % change to \reviewersinbindingfalse for submission

%------------------------------------------------
% Title page (no reviewers)
%------------------------------------------------
\begin{titlepage}
\centering
    {\Large \textbf{Freie Universität Berlin}}\\[1.5cm]
    {\large Institute of Meteorology}\\[0.5cm]
    {\large Department of Earth Sciences}\\[2.0cm]
    {\huge \textbf{Atmospheric Blocking and Warm Conveyor Belts:
    Mathematical Perspectives}}\\[1.5cm]
    {\large A cumulative dissertation submitted in partial fulfillment of the
    requirements for the degree of}\\[0.5cm]
    {\Large Doctor rerum naturalium (Dr.~rer.~nat.)}\\[1.5cm]
    {\large by}\\[0.5cm]
    {\Large \textbf{Henry Leon Schoeller}}\\[2.0cm]
\vfill
    {\large Berlin, June 2026}\\
\end{titlepage}
\cleardoublepage

%------------------------------------------------
% Reviewer page (separate unbound sheet at submission;
% bound into thesis in final version)
%------------------------------------------------
\ifreviewersinbinding
    % Final version: included in pagination as normal
    \chapter*{Examination Committee}
    \addcontentsline{toc}{chapter}{Examination Committee}
    \begin{center}
    \vspace*{3cm}

\else
    % Submission version: excluded from pagination
    \clearpage
    \thispagestyle{empty}
    \pagenumbering{gobble}
    \begin{center}
    \vspace*{3cm}

    {\Large \textbf{Examination Committee}}\\[2cm]
\fi


\begin{tabular}{rl}
    Chair:              & Prof.\ Dr.\ Fabian Hoffmann \\
    Deputy Chair:       & Prof.\ Dr.\ Henning Rust \\
    Committee Member:   & Dr.\ Mona Bukenberger \\[1em]
    Supervisor:         & Prof.\ Dr.\ Stephan Pfahl \\
    Reviewer:           & Prof.\ Dr.\ Gabriele Messori \\[1em]
    Disputation date:   & 14\textsuperscript{th} September 2026 \\
\end{tabular}
\end{center}
\ifreviewersinbinding\else
    % Restore pagination after the separate sheet
    \clearpage
    \pagenumbering{roman}
\fi
\cleardoublepage

%------------------------------------------------
% Declaration / Statement of originality
%------------------------------------------------
\chapter*{Declaration}
\addcontentsline{toc}{chapter}{Declaration}
I hereby declare that this dissertation has been composed by myself and that
the work contained herein is my own, except where explicitly stated otherwise.
This dissertation has not been submitted, in whole or in part, for any other
degree or professional qualification.

\medskip
\noindent
This cumulative dissertation is based on the following peer-reviewed
publications, to each of which I have contributed as detailed at the
end of the respective chapter \textbf{Needs Updates}:

\begin{enumerate}
    \item \fullcite{schoellerNoiseinducedEnhancementRegime2026}
    \item \fullcite{schoellerIsolatingEffectsModel2025}
    \item \fullcite{schoellerAssessingLagrangianCoherence2025}
\end{enumerate}

\vspace{1.5cm}
\noindent
Berlin, June 2026\\[1.5cm]
\noindent
\textbf{Henry Leon Schoeller}
\cleardoublepage

%------------------------------------------------
% Abstract
%------------------------------------------------
\chapter*{Abstract}
\addcontentsline{toc}{chapter}{Abstract}

Atmospheric blocking exerts a major influence on mid-latitude circulation
and is frequently associated with extreme weather, yet its
formation, persistence, and representation remain incompletely understood.
This is because atmospheric blocking is a complex phenomenon whose constitutive
processes bridge vast ranges of spatial and temporal scales -- prominently
through the interaction with warm conveyor belts. This work demonstrates that
viewing blocking, and in particular this interaction, through the lens of
dynamical systems theory is exceptionally well suited to capturing its
multi-scale nature, with regime persistence, coherence
under perturbations and state-dependent uncertainty as key concepts. 

The influence of small-scale processes on the large-scale flow is often
conceptualized as small random perturbations. Such perturbations can have a
first-order effect on the behaviour of a dynamical system and its stability
properties. Employing a paradigm model of mid-latitude atmospheric dynamics,
it is shown that noise can increase the lifetime of blockings and a method is
developed that can quantify this effect without access to the underlying
dynamics. Since the representation of small-scale processes is known to affect
blockings as well as warm conveyor belts, it is argued that this
mechanism's understanding may contribute to improving the representation of
their interaction.

Apart from errors in the representation of small-scale processes, uncertainty
in the estimated state of the atmosphere is the second major source of deficient
understanding and forecasting. By isolating short- from long-time scale
uncertainty components, the state-dependence of this uncertainty is studied
climatologically leveraging the ERA5 reanalysis. This provides coherent patterns
of anomalous
uncertainty for a set of Euro-Atlantic weather regimes -- including blockings --
as well as warm conveyor belts that can be explained by physical and mathematical
arguments. Moisture presence is a key contributor to elevated observation
uncertainties in addition to its significant role in model deficiencies.

While the preceding considerations concern how perturbations and uncertainty
affect the representation of blocking, its persistence ultimately rests on the
air streams that constitute it. Blocks are three-dimensional structures whose
properties are determined to a large degree by the air they consist of. The
origin of this air is systematically investigated by tracking Lagrangian
trajectories backward in time and clustering them according to their robustness
against small random perturbations. This conceptualization as coherent
trajectory bundles provides
 important insights into the kinematics of blockings and the surrounding flow 
and warm conveyor belts are identified among them. Reflecting on warm conveyor belts'
nature as spatio-temporally extended features is vital for understanding their
interaction with blockings.
\cleardoublepage

%------------------------------------------------
% Table of contents and lists
%------------------------------------------------
\tableofcontents
\cleardoublepage